%
% vibe-exit-code.tex
%
% Z Specification for the redesigned vox auto-vibe mood state machine
% (bead vox-ek1m). The redesign REPLACES an output-pattern classifier: the
% session's TTS mood is now derived from command EXIT CODES only -- no output
% text parsing, no command regex. A single Bash command yields ok (exit 0) or
% fail (non-zero); the mood is a total function of a bounded window of the most
% recent outcomes.
%
% This is the DESIGN-phase formal model. It must be fuzz-clean before the
% implementation mission dispatches. Section "Findings from Formalisation"
% records what the model reveals about the threshold boundaries, the K/run
% interaction, the relieved->happy decay, the empty window, and MAX-eviction.
%

\documentclass[a4paper,10pt,fleqn]{article}
\usepackage[margin=1in]{geometry}
\usepackage{fuzz}
\usepackage[colorlinks=true,linkcolor=blue,citecolor=blue,urlcolor=blue]{hyperref}

\begin{document}

\title{Vox Auto-Vibe from Exit Codes: A Z Specification}
\author{Formal model of the mood state machine (vox-ek1m)}
\date{July 2026}
\hypersetup{
  pdftitle={Vox Auto-Vibe from Exit Codes: A Z Specification},
  pdfauthor={Punt Labs},
  pdfsubject={Z Specification},
  pdfcreator={fuzz/probcli}
}
\maketitle

\tableofcontents
\newpage

\section{Introduction}

Auto-vibe chooses the voice the session speaks in. The old design read the
\emph{text} a command produced and matched it against patterns; the redesign
reads only the command's \emph{exit code}. Every Bash command contributes one
\emph{outcome} --- \texttt{ok} for exit~0, \texttt{fail} for any non-zero exit
--- and the session's mood is a total function of a bounded \emph{window} of the
most recent outcomes, newest last.

The design goal is an emotional arc that falls out of the numbers rather than
being asserted case by case:
\[
  happy \xrightarrow{\ fail\ } focused \xrightarrow{\ more\ } frustrated
  \xrightarrow{\ sustained\ } weary,
  \qquad
  \cdots \xrightarrow{\ ok\ } relieved \xrightarrow{\ clean\ } happy.
\]
Two numbers drive it. \emph{run} is the length of the trailing run of
consecutive failures at the end of the window; it measures how deep the current
trouble is. \emph{recentFail} asks only whether a failure appears among the last
$K$ outcomes; it is what separates a freshly-recovered session (\texttt{relieved})
from one that has been clean for a while (\texttt{happy}). The whole point of the
redesign is that an empty or clean window is \texttt{happy} --- the session never
\emph{defaults} to \texttt{frustrated}.

The model type-checks under \texttt{fuzz -t}. Everything the informal design
leaves to intuition --- where the thresholds sit, how $K$ interacts with the run
thresholds, exactly when relief decays to happiness, what an almost-empty window
does, and whether discarding the oldest outcome can ever change the mood ---
becomes a proof obligation here. Section~\ref{sec:findings} collects the answers.

\section{Basic Types}

Two free types carry the whole model. An \texttt{Outcome} is the classification
of a single command by its exit code. A \texttt{Mood} is one of the five voices
auto-vibe can select.

\begin{zed}
  Outcome ::= ok | fail
\end{zed}

\begin{zed}
  Mood ::= happy | focused | frustrated | weary | relieved
\end{zed}

The Boolean type: \texttt{fuzz} has none built in, and \texttt{ZBOOL} avoids the
\texttt{BOOL}/\texttt{true}/\texttt{false} keyword clash that would defeat ProB.

\begin{zed}
  ZBOOL ::= ztrue | zfalse
\end{zed}

\section{Constants}

Five constants fix the shape of the machine. \texttt{maxWindow} bounds the
window so the carrier is finite and ProB can animate it. \texttt{recentK} is the
width of the recency test that distinguishes \texttt{relieved} from
\texttt{happy}. The three \texttt{$\ldots$From} constants are the run thresholds:
the \emph{smallest} run that reads as \texttt{focused}, as \texttt{frustrated},
and as \texttt{weary} respectively.

The ordering constraints are not decoration. \texttt{focusFrom} $<$
\texttt{frustFrom} $<$ \texttt{wearyFrom} makes the three run-bands non-empty and
correctly ordered; \texttt{wearyFrom} $<$ \texttt{maxWindow} is what guarantees
the window is always long enough to \emph{reach} the deepest mood, so that
FIFO-eviction can never mask it (Finding~\ref{find:evict}). That
\texttt{focusFrom} is a \texttt{$\nat_1$} --- at least $1$ --- is what keeps the
\texttt{run}~$=0$ cases disjoint from the \texttt{focused} band
(Finding~\ref{find:focusone}).

\begin{axdef}
  maxWindow : \nat_1 \\
  recentK : \nat_1 \\
  focusFrom : \nat_1 \\
  frustFrom : \nat_1 \\
  wearyFrom : \nat_1
\where
  maxWindow = 20 \\
  recentK = 3 \\
  focusFrom = 1 \\
  frustFrom = 3 \\
  wearyFrom = 5 \\
  focusFrom < frustFrom \\
  frustFrom < wearyFrom \\
  wearyFrom < maxWindow
\end{axdef}

\section{Derivation}

The mood is derived, never stored. Two auxiliary functions capture the two
numbers, and \texttt{mood} composes them into a voice.

\subsection{The trailing run of failures}

\texttt{runFail}~$w$ is the length of the longest suffix of $w$ consisting
entirely of \texttt{fail}. It is characterised without recursion and without an
appeal to \texttt{max}: \texttt{runFail}~$w = n$ exactly when the last $n$
elements are all \texttt{fail} and the run stops there --- either $n$ exhausts
the window or the element just before the suffix is an \texttt{ok}. Exactly one
$n$ satisfies this for each $w$, so \texttt{runFail} is total and single-valued;
for the empty window the only witness is $n = 0$, giving
\texttt{runFail}~$\langle\rangle = 0$.

\begin{axdef}
  runFail : \seq Outcome \fun \nat
\where
  \forall w : \seq Outcome; n : \nat @ \\
  \quad runFail~w = n \iff \\
  \quad\quad (n \leq \# w \land \\
  \quad\quad~ (\forall i : (\# w - n + 1) \upto \# w @ w~i = fail) \land \\
  \quad\quad~ (n < \# w \implies w~(\# w - n) = ok))
\end{axdef}

\subsection{A failure in recent memory}

\texttt{recentFail}~$w$ holds exactly when some \texttt{fail} sits among the last
\texttt{recentK} positions of $w$. When the window is shorter than
\texttt{recentK} every position counts, so a short window with any failure at all
is ``recent''. An empty window has no such position, so \texttt{recentFail} is
\texttt{zfalse}.

\begin{axdef}
  recentFail : \seq Outcome \fun ZBOOL
\where
  \forall w : \seq Outcome @ \\
  \quad recentFail~w = ztrue \iff \\
  \quad\quad (\exists i : \dom w | i > \# w - recentK @ w~i = fail)
\end{axdef}

\subsection{The mood}

\texttt{mood} maps a window to a voice. The five clauses partition every window:
\texttt{run}~$=0$ splits on \texttt{recentFail} into \texttt{happy} and
\texttt{relieved}; a positive \texttt{run} falls into exactly one of the three
run-bands. The clauses are mutually exclusive and jointly exhaustive, so
\texttt{mood} is total and single-valued --- the constant \texttt{focusFrom}~$=1$
is what closes the gap that would otherwise sit between \texttt{run}~$=0$ and the
\texttt{focused} band.

\begin{axdef}
  mood : \seq Outcome \fun Mood
\where
  \forall w : \seq Outcome @ \\
  \quad (runFail~w = 0 \land recentFail~w = zfalse \implies mood~w = happy) \land \\
  \quad (runFail~w = 0 \land recentFail~w = ztrue \implies mood~w = relieved) \land \\
  \quad (runFail~w \in focusFrom \upto frustFrom - 1 \implies mood~w = focused) \land \\
  \quad (runFail~w \in frustFrom \upto wearyFrom - 1 \implies mood~w = frustrated) \land \\
  \quad (runFail~w \geq wearyFrom \implies mood~w = weary)
\end{axdef}

With the constants above the table reads: \texttt{run}~$=0$ gives \texttt{happy}
or \texttt{relieved}; \texttt{run} $\in 1 \upto 2$ gives \texttt{focused};
\texttt{run} $\in 3 \upto 4$ gives \texttt{frustrated}; \texttt{run} $\geq 5$
gives \texttt{weary}.

\section{State}

The state is a single component: the bounded window of outcomes, newest last.
Because the state carries exactly one variable, framing an operation is trivial
--- constraining \texttt{window'} constrains the whole state. The mood is not a
component; it is \texttt{mood}~\texttt{window}, recomputed on demand.

\begin{schema}{Vibe}
  window : \seq Outcome
\where
  \# window \leq maxWindow
\end{schema}

\section{Initialisation}

A fresh session starts with an empty window, and therefore --- by the derivation,
not by a special case --- with mood \texttt{happy}.

\begin{schema}{Init}
  Vibe'
\where
  window' = \langle \rangle
\end{schema}

The empty-window default is a consequence, not a stipulation. This property
schema restates it as a type-checked predicate; its truth is the second clause of
\texttt{mood} applied to $\langle\rangle$, for which \texttt{runFail} is $0$ and
\texttt{recentFail} is \texttt{zfalse}.

\begin{schema}{DefaultHappy}
  Vibe
\where
  window = \langle \rangle \implies mood~window = happy
\end{schema}

\section{Operations}

Each command records one outcome. The window is a FIFO buffer of capacity
\texttt{maxWindow}: appending to a full window evicts the oldest outcome.
\texttt{pushBounded} captures the append-with-eviction once, so the two record
operations differ only in the outcome they append.

\begin{axdef}
  pushBounded : (\seq Outcome \cross Outcome) \fun \seq Outcome
\where
  \forall w : \seq Outcome; o : Outcome @ \\
  \quad (\# w < maxWindow \implies pushBounded(w, o) = w \cat \langle o \rangle) \land \\
  \quad (\# w = maxWindow \implies pushBounded(w, o) = (tail~w) \cat \langle o \rangle)
\end{axdef}

\subsection{RecordOk}

A command exited $0$. Its outcome is appended; if the window was full the oldest
outcome is evicted.

\begin{schema}{RecordOk}
  \Delta Vibe
\where
  window' = pushBounded(window, ok)
\end{schema}

\subsection{RecordFail}

A command exited non-zero. Symmetric to \texttt{RecordOk}, appending
\texttt{fail}.

\begin{schema}{RecordFail}
  \Delta Vibe
\where
  window' = pushBounded(window, fail)
\end{schema}

\subsection{Observing the mood}

A pure query: the mood the session is currently in, with no change of state.

\begin{schema}{CurrentMood}
  \Xi Vibe \\
  mood! : Mood
\where
  mood! = mood~window
\end{schema}

\section{The emotional arc as properties}

The arc is not asserted; it is a consequence of the derivation and the FIFO
append. The following property schemas restate each leg as a type-checked
predicate over an operation. Their truth is argued in the prose that precedes
each; \texttt{fuzz} checks that they are well-typed, and small ProB animations
(recommended before implementation) check that they hold across the reachable
state space.

\subsection{A failure deepens, never eases}

Appending a \texttt{fail} lengthens the trailing run by one, capped at the window
size. In the operative range --- \texttt{run} $<$ \texttt{maxWindow}, which holds
whenever the mood is below \texttt{weary}-of-a-full-window --- the run simply
increments, so each additional failure can only push the mood \emph{deeper} into
the arc: \texttt{happy}/\texttt{relieved} $\to$ \texttt{focused} $\to$
\texttt{frustrated} $\to$ \texttt{weary}.

\begin{schema}{FailDeepens}
  RecordFail
\where
  runFail~window' \geq runFail~window \\
  runFail~window' \leq runFail~window + 1
\end{schema}

The first failure out of a clean or empty window makes \texttt{run}~$=1$, which
is \texttt{focused} --- the first step of the arc.

\begin{schema}{FailFromClean}
  RecordFail
\where
  runFail~window = 0 \implies mood~window' = focused
\end{schema}

\subsection{A success resets the run}

Appending an \texttt{ok} makes the last element \texttt{ok}, so the trailing run
of failures is $0$. A success therefore always lands the session in
\texttt{happy} or \texttt{relieved} and never leaves it \texttt{focused},
\texttt{frustrated}, or \texttt{weary}: the recovery side of the arc is immediate.

\begin{schema}{OkResetsRun}
  RecordOk
\where
  runFail~window' = 0
\end{schema}

\begin{schema}{OkNeverNegative}
  RecordOk
\where
  mood~window' = happy \lor mood~window' = relieved
\end{schema}

Which of the two it is depends solely on recency: a success immediately after
failures still has a failure within the last \texttt{recentK} positions, so it is
\texttt{relieved}; once enough successes have pushed every failure past that
window, \texttt{recentFail} turns \texttt{zfalse} and the mood decays to
\texttt{happy}.

\begin{schema}{OkRelievedIffRecent}
  RecordOk
\where
  mood~window' = relieved \iff recentFail~window' = ztrue
\end{schema}

\section{Findings from Formalisation}
\label{sec:findings}

Writing the model down settled six questions the informal design left open. Each
finding names the issue, states what the model shows, and where relevant
recommends a concrete constant.

\subsection{The empty window is happy by derivation, not by a special case}
\label{find:empty}

The informal spec lists ``window empty $\to$ happy'' as a rule of its own. The
model shows it is redundant: for $\langle\rangle$, \texttt{runFail} is $0$ (the
candidate set is $\{0\}$) and \texttt{recentFail} is \texttt{zfalse} (no position
satisfies the existential), so the second \texttt{mood} clause already yields
\texttt{happy}. \texttt{DefaultHappy} is a theorem, not an axiom. This is the
strongest form of ``never default to frustrated'': there is no default at all,
only the derivation, and the derivation of the empty window is \texttt{happy}.
\textbf{Recommendation:} do not implement a special empty-window branch; compute
\texttt{run} and \texttt{recentFail} uniformly and let the table decide.

\subsection{The threshold boundaries are unambiguous but must be stated as
half-open bands}
\label{find:bands}

The prose gives the bands as ``\texttt{run} in $1..2$'', ``$3..4$'', ``$\geq 5$''.
The model shows these are exactly the half-open intervals
$[\texttt{focusFrom}, \texttt{frustFrom})$, $[\texttt{frustFrom},
\texttt{wearyFrom})$, $[\texttt{wearyFrom}, \infty)$ with the current constants
$(1,3,5)$. The boundary cases are therefore fixed and worth stating in the
implementation's own words: \texttt{run}~$=2$ is the last \texttt{focused} value
and \texttt{run}~$=3$ the first \texttt{frustrated}; \texttt{run}~$=4$ is the last
\texttt{frustrated} and \texttt{run}~$=5$ the first \texttt{weary}. The bands are
correct \emph{iff} \texttt{focusFrom} $<$ \texttt{frustFrom} $<$
\texttt{wearyFrom}, which the constant schema asserts. \textbf{Recommendation:}
keep $(1,3,5)$. Two consecutive failures reading as \texttt{focused} (not yet
\texttt{frustrated}) matches how a developer feels about a flaky first retry; five
in a row reading as \texttt{weary} is a defensible ``this is not going well''
signal. If tuning is wanted later, change only these three constants --- the
derivation is parameterised on them and nothing else moves.

\subsection{\texttt{focusFrom} must be exactly $1$, or totality breaks}
\label{find:focusone}

The single most important thing the formalisation caught. If \texttt{focusFrom}
were $0$, the \texttt{run}~$=0$ window would satisfy \emph{both} the
\texttt{happy}/\texttt{relieved} clauses and the \texttt{focused} clause ---
\texttt{mood} would be over-determined and inconsistent. If \texttt{focusFrom}
were $\geq 2$, runs in $1 \upto \texttt{focusFrom} - 1$ would satisfy \emph{no}
clause --- \texttt{mood} would be underspecified there and not total. Only
\texttt{focusFrom}~$=1$ makes the five clauses a clean partition. The
\texttt{$\nat_1$} declaration rules out $0$; the equality $= 1$ and the disjoint
bands rule out $\geq 2$. \textbf{Recommendation:} the implementation must treat
``\texttt{run}~$\geq 1$'' as the entry to \texttt{focused} --- there is no free
parameter here; $1$ is forced.

\subsection{\texttt{recentK} interacts with the run thresholds only in the
\texttt{run}~$=0$ case, and \texttt{relieved} is unambiguous}
\label{find:kinteract}

The question was whether a window could be both ``\texttt{run}~$=0$,
\texttt{recentFail}'' and something else. It cannot. Whenever
\texttt{run}~$\geq 1$ the last outcome is a \texttt{fail}, so \texttt{recentFail}
is necessarily \texttt{ztrue} (for any \texttt{recentK}~$\geq 1$) --- but the mood
in that case is decided by the run-band, and \texttt{recentFail} is simply not
consulted. \texttt{recentFail} changes the answer \emph{only} when
\texttt{run}~$=0$, precisely the \texttt{happy}-vs-\texttt{relieved} split. So
\texttt{relieved} is unambiguous: it is exactly ``the last outcome is
\texttt{ok}, but a \texttt{fail} occurred within the last \texttt{recentK}''.
\textbf{Recommendation:} keep \texttt{recentK}~$=3$. Note the constraint that
matters is \texttt{recentK}~$\geq 1$; the specific value $3$ tunes only how many
clean commands it takes to decay from \texttt{relieved} to \texttt{happy} (see
below), and does not affect any other mood.

\subsection{Relief decays to happy after exactly \texttt{recentK} clean commands
--- provided no eviction intervenes}
\label{find:decay}

``Exactly when does \texttt{recentFail} become false?'' The model gives a precise
answer: \texttt{recentFail}~$w$ is \texttt{zfalse} iff none of the last
\texttt{recentK} positions is a \texttt{fail}. Starting from a window whose last
outcome is a failure and appending only successes, the failure is at distance $1$
from the end after the first \texttt{ok}, distance $2$ after the second, and
distance \texttt{recentK} after the \texttt{recentK}-th; on the outcome that
pushes it to distance \texttt{recentK}~$+1$ it leaves the recency window and the
mood decays \texttt{relieved}~$\to$~\texttt{happy}. With \texttt{recentK}~$=3$
that is the third clean command after the last failure --- assuming the window is
long enough to still hold that failure. This is the one real subtlety: if the
last failure has already been evicted (below), \texttt{recentFail} can read
\texttt{zfalse} sooner, which only \emph{hastens} the decay to \texttt{happy} ---
a benign direction. \textbf{Recommendation:} implement \texttt{recentFail} as
``any of the last \texttt{recentK} recorded outcomes is a failure'', reading from
the stored window, so decay timing follows the window automatically.

\subsection{FIFO-eviction can change the derived mood only in the benign
direction, and never at all below \texttt{weary}}
\label{find:evict}

The worry: can dropping the oldest outcome flip the mood surprisingly? The model
bounds it exactly. Eviction removes \texttt{window}'s \emph{first} element.
\texttt{runFail} depends only on the trailing suffix, so eviction changes
\texttt{runFail} only when the entire window is failures --- \texttt{run}~$=
\texttt{maxWindow}$ --- in which case appending one more failure and evicting the
oldest leaves \texttt{run} at \texttt{maxWindow} rather than growing it. Because
\texttt{wearyFrom}~$<$~\texttt{maxWindow} ($5 < 20$), the window is always long
enough to represent every band up to and including \texttt{weary} without
eviction touching the run; eviction can cap the run at $20$, but $20$ and $5$ both
read as \texttt{weary}, so the \emph{mood} is unchanged. \texttt{recentFail} is
similarly affected only if the failure it was counting is the one evicted, and
that can only turn \texttt{recentFail} off --- moving \texttt{relieved} toward
\texttt{happy}, the recovery direction. There is no transition in which eviction
makes the mood \emph{worse} (deeper into the arc) or produces a mood outside the
table. \textbf{Recommendation:} keep \texttt{maxWindow}~$=20$. Any value $>$
\texttt{wearyFrom} is safe; $20$ gives comfortable headroom and keeps the
carrier small enough for ProB to animate. The load-bearing invariant to preserve
in code is \texttt{wearyFrom}~$<$~\texttt{maxWindow}.

\subsection{Summary of recommended constants}

\begin{center}
\begin{tabular}{lll}
\textbf{Constant} & \textbf{Value} & \textbf{Constraint the model requires} \\
\hline
\texttt{focusFrom}  & $1$  & forced: totality demands exactly $1$ \\
\texttt{frustFrom}  & $3$  & \texttt{focusFrom} $<$ \texttt{frustFrom} \\
\texttt{wearyFrom}  & $5$  & \texttt{frustFrom} $<$ \texttt{wearyFrom} \\
\texttt{recentK}    & $3$  & $\geq 1$; tunes decay speed only \\
\texttt{maxWindow}  & $20$ & $>$ \texttt{wearyFrom}; keeps carrier finite \\
\end{tabular}
\end{center}

\end{document}
